\section{Functions}

\begin{definition} \label{an:def:function}
  Suppose a relation $G$ from $X$ to $Y$ satisfies
  \[
    \forall x \in X, \exists! y \in Y, (x, y) \in G
  \]
  Then the ordered triple
  \[
    f \coloneqq (X, Y, G)
  \]
  is called a \textbf{function}, where
  \begin{itemize}
    \item $X$ is called the \textbf{domain} of $f$
    \item $Y$ is called the \textbf{codomain} of $f$
    \item $G$ is called the \textbf{graph} of $f$
  \end{itemize}
  The notation
  \begin{itemize}
    \item $f:X \to Y$ denotes that $f$ is a function from $X$ to $Y$
    \item $f(x) = y$ denotes $(x, y)\in G$
  \end{itemize}
\end{definition}

\bigskip

\begin{definition} \label{an:def:image-inverse-image}
  Suppose $f: X \to Y$.
  \begin{itemize}
    \item The \textbf{image} of $E \subseteq X$ under $f$ is defined by
          \[
            f(E) \coloneqq \setbuilder*{f(x)}{x \in E} \subseteq Y
          \]
    \item The \textbf{inverse image} of $E \subseteq Y$ under $f$ is defined by
          \[
            f^{-1}(E) \coloneqq \setbuilder*{x \in X}{f(x) \in E} \subseteq X
          \]
  \end{itemize}
\end{definition}

\bigskip

\begin{definition} \label{an:def:injective-surjective-bijective}
  Suppose $f: X \to Y$.
  \begin{itemize}
    \item $f$ is said to be \textbf{injective} if
          \[
            \forall x_{1}, x_{2} \in X, f(x_{1}) = f(x_{2}) \implies x_{1} = x_{2}
          \]
          The notation $f: X \hookrightarrow Y$ denotes that $f$ is injective.
    \item $f$ is said to be \textbf{surjective} if
          \[
            \forall y \in Y, \exists x \in X, f(x) = y
          \]
          or equivalently,
          \[
            f(X) = Y
          \]
          The notation $f: X \twoheadrightarrow Y$ denotes that $f$ is surjective.
    \item $f$ is said to be \textbf{bijective} if
          \begin{itemize}
            \item $f$ is injective
            \item $f$ is surjective
          \end{itemize}
          or equivalently,
          \[
            \forall y \in Y, \exists! x \in X, f(x) = y
          \]
          The notation $f: X \xrightarrow{\sim} Y$ denotes that $f$ is bijective.
  \end{itemize}
\end{definition}

\bigskip

\begin{definition} \label{an:def:inverse-function}
  Suppose $f: X \xrightarrow{\sim} Y$. The \textbf{inverse function} $f^{-1}: Y \to X$ of $f$ is defined by
  \[
    \forall (y, x) \in Y \times X, f^{-1}(y) = x \iff f(x) = y
  \]
  By \cref{an:def:injective-surjective-bijective},
  \[
    \begin{array}{c}
      \forall y \in Y, \exists! x \in X, f(x) = y \\
      \Downarrow                                  \\
      \forall y \in Y, \exists! x \in X, f^{-1}(y) = x
    \end{array}
  \]
  Therefore, $f^{-1}$ is well-defined.
\end{definition}

\bigskip

\begin{theorem} \label{an:thm:inverse-function-is-bijective}
  Suppose $f: X \xrightarrow{\sim} Y$. Then
  \[
    f^{-1}: Y \xrightarrow{\sim} X
  \]
\end{theorem}

\begin{proof}
  By \cref{an:def:function},
  \[
    \begin{array}{c}
      \forall x \in X, \exists! y \in Y, f(x) = y \\
      \Downarrow                                  \\
      \forall x \in X, \exists! y \in Y, f^{-1}(y) = x
    \end{array}
  \]
\end{proof}

\bigskip

\begin{definition} \label{an:def:restriction-function}
  Suppose
  \begin{itemize}
    \item $f: X \to Y$
    \item $E \subseteq X$
    \item $F \subseteq Y$
  \end{itemize}
  Then
  \begin{itemize}
    \item The \textbf{restriction} $f|_{E}: E \to Y$ of $f$ to $E$ is defined by
          \[
            \forall x \in E, f|_{E}(x) = f(x)
          \]
    \item The \textbf{corestriction} $f|^{F}: X \to F$ of $f$ to $F \supseteq f(X)$ is defined by
          \[
            \forall x \in X, f|^{F}(x) = f(x)
          \]
    \item The \textbf{bi-restriction} $f|_{E}^{F}: E \to F$ of $f$ to $E$ and $F \supseteq f(E)$ is defined by
          \[
            \forall x \in E, f|_{E}^{F}(x) = f(x)
          \]
  \end{itemize}
\end{definition}

\bigskip

\begin{theorem} \label{an:thm:properties-of-restriction-functions}
  Suppose
  \begin{itemize}
    \item $f: X \to Y$
    \item $E \subseteq X$
    \item $f(E) \subseteq F \subseteq Y$
  \end{itemize}
  Then
  \begin{enumerate}[label=\textbf{(\Alph*)}]
    \item $f: X \hookrightarrow Y \implies f|_{E}^{F}: E \hookrightarrow F$
    \item $f(E) = F \implies f|_{E}^{F}: E \twoheadrightarrow F$
    \item $f: X \hookrightarrow Y \land f(E) = F \implies f|_{E}^{F}: E \xrightarrow{\sim} F$
  \end{enumerate}
\end{theorem}

\begin{proof}
  \leavevmode\par
  \begin{enumerate}[label=\textbf{(\Alph*)}]
    \item Suppose
          \begin{itemize}
            \item $f: X \hookrightarrow Y$
            \item $x_{1}, x_{2} \in E$
            \item  $f|_{E}^{F}(x_{1}) = f|_{E}^{F}(x_{2})$
          \end{itemize}
          By \cref{an:def:injective-surjective-bijective},
          \[
            \begin{array}{c}
              f(x_{1}) = f(x_{2}) \\
              \Downarrow          \\
              x_{1} = x_{2}
            \end{array}
          \]
    \item Suppose $f(E) = F$. By \cref{an:def:image-inverse-image},
          \[
            \begin{array}{c}
              \forall y \in F, \exists x \in E, f(x) = y \\
              \Downarrow                                 \\
              \forall y \in F, \exists x \in E, f|_{E}^{F}(x) = y
            \end{array}
          \]
    \item follows from \textbf{(A)} and \textbf{(B)}.
  \end{enumerate}
\end{proof}

\bigskip

\begin{definition} \label{an:def:composition-function}
  Suppose
  \begin{itemize}
    \item $f: X \to Y$
    \item $g: Y \to Z$
  \end{itemize}
  The \textbf{composition} $g \circ f: X \to Z$ of $g$ with $f$ is defined by
  \[
    \forall x \in X, (g \circ f)(x) = g(f(x))
  \]
\end{definition}

\bigskip

\begin{theorem} \label{an:thm:properties-of-composition-functions}
  Suppose
  \begin{itemize}
    \item $f: X \to Y$
    \item $g: Y \to Z$
  \end{itemize}
  Then
  \begin{enumerate}[label=\textbf{(\Alph*)}]
    \item $f: X \hookrightarrow Y \land g: Y \hookrightarrow Z$ $\implies g \circ f: X \hookrightarrow Z$
    \item $f: X \twoheadrightarrow Y \land g: Y \twoheadrightarrow Z$ $\implies g \circ f: X \twoheadrightarrow Z$
    \item $f: X \xrightarrow{\sim} Y \land g: Y \xrightarrow{\sim} Z$ $\implies g \circ f: X \xrightarrow{\sim} Z$
  \end{enumerate}
\end{theorem}

\begin{proof}
  \leavevmode\par
  \begin{enumerate}[label=\textbf{(\Alph*)}]
    \item Suppose
          \begin{itemize}
            \item $f: X \hookrightarrow Y$
            \item $g: Y \hookrightarrow Z$
            \item $x_{1}, x_{2} \in X$
            \item $(g \circ f)(x_{1}) = (g \circ f)(x_{2})$
          \end{itemize}
          By \cref{an:def:injective-surjective-bijective},
          \[
            \begin{array}{c}
              f(x_{1}) = f(x_{2}) \\
              \Downarrow          \\
              x_{1} = x_{2}
            \end{array}
          \]
    \item Suppose
          \begin{itemize}
            \item $f: X \twoheadrightarrow Y$
            \item $g: Y \twoheadrightarrow Z$
          \end{itemize}
          By \cref{an:def:injective-surjective-bijective},
          \begin{itemize}
            \item $\forall z \in Z, \exists y \in Y, g(y) = z$
            \item $\forall y \in Y, \exists x \in X, f(x) = y$
          \end{itemize}
          Therefore,
          \[
            \forall z \in Z, \exists x \in X, (g \circ f)(x) = z
          \]
    \item follows from \textbf{(A)} and \textbf{(B)}.
  \end{enumerate}
\end{proof}
